% !TEX root = ../main.tex

\section{Method}
\label{sec:method}

\subsection{Problem Setting}

Tasks arrive in a time-ordered stream,
\[
x_1, x_2, \ldots, x_T, \qquad x_t \sim P_t,
\]
where $P_t$ is the (potentially shifting) task distribution at time $t$ and each $x_t$ has a fixed ground truth $y^*(x_t)$. The agent observes the history of prior tasks, their trajectories, and outcomes,
\[
\Hist = \{(x_i, r_i, \tau_i)\}_{i=1}^{t-1}.
\]
A \emph{harness} $C = \varphi(\Hist)$ with $|C| \leq K$ is a bounded representation (prompts, skills, memory, tools, infrastructure) with capacity budget $K$, and the agent acts via $a_t \sim \pi(a \mid x_t, C_t)$. Policy stochasticity (sampling, reasoning) is the only randomness in the system; a better harness concentrates the policy on correct actions.

\subsection{Conceptual Framework}
\label{sec:framework}

\paragraph{Utility and regret.}
The utility of a harness $C$ on a task $x$ is $V(C, x) = \mathbb{E}_{a \sim \pi(\cdot \mid x, C)}[r(a, y^*(x))]$, and the regret of the evolver-constructed harness relative to the full-history policy is
\[
\text{Regret}(\varphi, x_t) = V(\Hist, x_t) - V(\varphi(\Hist), x_t) \geq 0,
\]
non-negative because $\Hist$ strictly contains more information than its bounded compression.

\paragraph{Two dimensions of harness construction.}
Fix the evolver class $\Phi$. The best task-conditional harness achievable within $\Phi$ at capacity~$K$ is
\[
C^*_\Phi(x) = \arg\max_{\varphi \in \Phi,\, |C| \leq K} V\!\bigl(\varphi(\Hist, x),\, x\bigr).
\]
Two dimensions matter: the construction quality of $\Phi$ (what harnesses it can build) and whether the harness is \emph{conditional} on $x_t$ (per-task routing) or \emph{unconditional} (one $C$ for all tasks).

\begin{proposition}[Regret Decomposition]
\label{prop:decomp}
For any unconditional evolver $\varphi \in \Phi$:
\[
\mathbb{E}_{x_t}[\text{Regret}(\varphi, x_t)] = \Levo(\Phi, K) + \Lnav(\varphi),
\]
where
\begin{align*}
\Levo(\Phi, K) &\!=\! \mathbb{E}_{x_t}\!\bigl[V(\Hist, x_t)\!-\! V(C^*_\Phi(x_t), x_t)\bigr], \\
\Lnav(\varphi)  &\!=\! \mathbb{E}_{x_t}\!\bigl[V(C^*_\Phi(x_t), x_t)\!-\! V(\varphi(\Hist), x_t)\bigr].
\end{align*}
Both terms are non-negative and the decomposition telescopes. $\Levo$ is the \emph{evolver-class limit}; $\Lnav$ is the \emph{conditioning limit}.
\end{proposition}

\begin{proposition}[Stationarity $\Rightarrow \Lnav{=}0$]
\label{prop:stationarity}
Under task-stream stationarity --- there exists a single $C^*$ optimal for every task in $\mathrm{support}(P_t)$ --- the optimal unconditional harness matches the per-task optimum and $\Lnav = 0$. This explains why linear-chain evolution succeeds on i.i.d.\ benchmarks.
\end{proposition}

\begin{proposition}[$\Lnav$ grows with shift]
\label{prop:shift}
Let shift magnitude $\Shift = \mathbb{E}_{x, x' \sim P_t}[d(x, x')]$, where, abbreviating $v(a,b) \!\equiv\! V(C^*_\Phi(a), b)$,
\[
d(x, x') = v(x, x) + v(x', x') - v(x, x') - v(x', x)
\]
measures how much two tasks disagree on which harness is better. Then $\Shift > 0 \Rightarrow \Lnav(\varphi^*_{\text{uncond}}) > 0$, and $\Lnav$ is monotone in $\Shift$: as tasks diverge, a single compromise harness loses more.
\end{proposition}

\begin{proposition}[Temporal accumulation]
\label{prop:accum}
Under non-stationarity the task distribution diversifies over time, so $\Shift$ grows with $t$ and $d\Lnav/dt > 0$. The total regret dynamic is therefore
\[
\tfrac{d}{dt}\mathbb{E}[\text{Regret}] = \tfrac{d\Levo}{dt} + \tfrac{d\Lnav}{dt},
\]
with $d\Levo/dt \leq 0$ (the evolver learns) and $d\Lnav/dt > 0$ (shift accumulates). Learning initially outpaces shift and accuracy improves; eventually shift dominates and accuracy degrades. The resulting non-monotonic curve is what makes early-freeze beat full evolution.
\end{proposition}

\paragraph{Predictive and diagnostic use of the decomposition.}
The two losses admit qualitatively different kinds of analysis. $\Lnav$ is \emph{empirically predictive}: $\Shift$ is measurable from per-batch task heterogeneity or from oracle-branch gap variance, and Propositions~\ref{prop:stationarity}--\ref{prop:accum} yield falsifiable claims that we test in~\S\ref{sec:navloss}. $\Levo$ is \emph{diagnostic vocabulary}: evolver-class capability has no natural scalar parameterization --- ``stronger evolver'' is not a continuous axis --- so direct quantification of $\Levo$ is not attempted. Instead the multi-agent contribution (\S\ref{sec:method-multi}) is specified as a set of architectural patterns that prior single-agent evolvers lack, and its value is witnessed by ablation and artifact-type analysis (\S\ref{sec:ablation}) rather than by a class-expansion bound.

\paragraph{Experience availability as a third axis.}
When $\Hist$ contains no relevant signal for the current task's regime --- missing credentials, novel language, unfamiliar data source --- neither expanding $\Phi$ (better construction given signal) nor conditioning $\varphi(\Hist, x)$ (better routing among existing signal) can help. This outside-the-decomposition regime motivates the third intervention, HITL~(\S\ref{sec:method-hitl}), which supplies direction at the evolver's \emph{input} rather than modifying either loss term. The decomposition therefore yields a regime-to-remedy mapping: stationary streams require nothing beyond a single evolver; non-stationary streams with simple harnesses are addressed by navigation; non-stationary streams with complex harnesses by multi-agent evolution and navigation jointly; and novel regimes outside $\Hist$ by HITL.

\begin{figure*}[t]
\centering
\rule{0.95\textwidth}{6.5cm}
\caption{Unified system pipeline. \emph{Top}: a git-backed strategy tree provides per-task routing --- $F_{\text{Navigate}}$ reads each branch's workspace (\texttt{README.md}, \texttt{prompts/system.md}, \texttt{skills/}, \texttt{tools/registry.yaml}) via \texttt{git show} and routes $x_t$ to the best-equipped branch (\S\ref{sec:method-nav}). \emph{Bottom}: per evolution cycle, a four-phase multi-agent pipeline operates on the routed branch --- Analyst $\to$ parallel Researchers $\to$ Builder $\to$ Verifier --- with persistent cross-cycle state (\texttt{task\_board.md}, \texttt{research\_log.jsonl}, \texttt{architecture.md}) linking successive cycles (\S\ref{sec:method-multi}). Two HITL hook points are marked: a \emph{taskboard direction} hook between Analyst and Researcher (Cycle~$k$, addresses C3) and a \emph{builder credential} hook inside Builder execution (Cycle~$k+m$, resource gate); both are silent unless triggered (\S\ref{sec:method-hitl}).}
\label{fig:unified_pipeline}
\end{figure*}

\subsection{Agentic Navigation}
\label{sec:method-nav}

Propositions~\ref{prop:shift} and~\ref{prop:accum} identify per-task conditioning as the principled response to $\Lnav$. We instantiate conditioning as a git-backed strategy tree with per-task routing (top of Figure~\ref{fig:unified_pipeline}). The solver's workspace is a git repository whose \texttt{main} branch carries stationary content (techniques that transfer across regimes), with regime-specific branches (e.g., \texttt{branch/crypto-classical}, \texttt{branch/binary-reversing}) spawned by the evolver when distribution shift is detected. Each branch maintains its own prompt, skills, and tool registry.

Before solving each task, a navigator LLM reads the workspace content of every branch --- \texttt{README.md}, \texttt{prompts/system.md}, the \texttt{skills/} directory listing, and \texttt{tools/registry.yaml} --- via \texttt{git show}, and emits \texttt{\{"branch": \ldots, "confidence": \ldots, "reason": \ldots\}} routing the task to the best-equipped branch. The solver checks out the selected branch and executes on that workspace. We ship two evolution modes: \emph{inline}, a single sandboxed evolver call with full git access that can create and mutate branches, and \emph{orchestrated}, a plan-driven variant with a separate analyst pass and per-branch evolver calls.

This design is distinct from retrieval-augmented generation and mixture-of-experts routing in four respects: retrieval units are executable git branches (workspaces), not document chunks; branches carry evolutionary lineage via fork and merge; the router operates over README and system-prompt content rather than embedding similarity; and total capacity is partitioned across branches rather than retrieved as top-$k$.

\subsection{Multi-Agent Evolution}
\label{sec:method-multi}

To address C2 the evolver must be able to construct multi-file infrastructure, not merely edit prompts. Rather than asserting a class-expansion theorem, we identify four architectural patterns that prior single-agent evolvers lack and whose joint presence supports the construction of infrastructure-class artifacts.

The first pattern is \emph{persistent cross-cycle state}: three files persist across evolution cycles --- \texttt{task\_board.md} (failure regimes and priority tags, updated by the Analyst), \texttt{research\_log.jsonl} (one JSON record per tested approach, including \texttt{works:true/false}, coverage, and latency), and \texttt{architecture.md} (what the Builder has already constructed). These let cycle $N$ build on cycle $N{-}1$'s findings, which single-cycle chain evolvers cannot.

The second pattern is \emph{parallel independent research}: in each cycle, multiple Researcher agents explore disjoint hypothesis branches in isolation (no shared conversation state), which reduces premature convergence compared to sequential refinement.

The third pattern is \emph{role-distinct objectives}: Analyst decomposes failures, Researchers investigate, Builder implements, and Verifier tests, with the Verifier using a success criterion separate from the Builder's --- unlike single-objective self-critique as used in Reflexion or Self-Refine.

The fourth pattern is \emph{HITL phase-boundary hooks}, described in~\S\ref{sec:method-hitl}.

The four-phase cycle is shown in the bottom of Figure~\ref{fig:unified_pipeline}. In Phase~1 (Analyze) the Analyst reads batch trajectories and updates \texttt{task\_board.md}. In Phase~2 (Research) $N$ parallel sandboxed Researchers with network access test data sources, APIs, and techniques per regime, appending to \texttt{research\_log.jsonl}. In Phase~3 (Build) the Builder reads only the \texttt{works:true} records, constructs \texttt{infra/*.py} pipelines, and updates \texttt{architecture.md}. In Phase~4 (Verify) the Verifier tests each new artifact against sample queries; failed artifacts feed back to the Builder for up to three retries.

\subsection{Human-in-the-Loop Integration}
\label{sec:method-hitl}

When $\Hist$ contains no useful signal for $x_t$'s regime the evolver's \emph{input} is the bottleneck, not its construction capability. Expanding~$\Phi$ or conditioning $\varphi$ does not repair this. We therefore add triggered hooks at two structurally distinct points in the multi-agent pipeline (marked in Figure~\ref{fig:unified_pipeline}), differentiated by \emph{what} the human supplies.

The first is a \emph{taskboard direction} hook between Analyst and Researcher: when the Analyst flags a regime with no covering branch in the strategy tree and no prior \texttt{works:true} signal in \texttt{research\_log.jsonl}, a one-line direction prompt is offered (e.g., ``describe how to approach \texttt{<regime>} or skip''). The human supplies \emph{semantic guidance} that shapes the entire research cycle. This hook addresses C3 directly: it injects a transferable prior when $\Hist$ has none.

The second is a \emph{builder credential} hook inside Phase~3: when the Builder's pipeline-construction step encounters an authentication failure (\texttt{401}/\texttt{403}) or a configuration prompt for missing credentials, a token-injection prompt is offered (e.g., ``provide \texttt{EASTMONEY\_API\_KEY} or skip''). The human supplies a \emph{resource} the evolver cannot discover from trajectories alone. This is mechanistically distinct from C3 (the data exists; access is gated), but shares the empirical signature that no amount of additional evolution closes the gap.

Both hooks are silent unless triggered; they fire on evolver-internal signals (Analyst-flagged novel regime; HTTP authentication failure during Builder execution) and otherwise impose zero cost. The typical invoked cost is under thirty seconds of human time per cycle.
